% concatenated from final.tex
\documentclass[10 pt, reqno]{amsart}

\usepackage{mathtools}
\DeclarePairedDelimiter\ceil{\lceil}{\rceil}
\DeclarePairedDelimiter\floor{\lfloor}{\rfloor}

\usepackage[dvipsnames]{xcolor}
\textheight 22.5truecm \textwidth 14.5truecm
\setlength{\oddsidemargin}{0.35in}\setlength{\evensidemargin}{0.35in}

\setlength{\topmargin}{-.5cm}



\usepackage[a4paper,margin=1in]{geometry}
\usepackage{latexsym, mathrsfs, enumerate, fullpage,graphicx,amsmath,amsfonts,amssymb,amsthm,tikz-cd}

% \doublespacing
\usepackage{biblatex}
\addbibresource{mybib.bib}
\usepackage{hyperref}
% main shared counter (reset per section)
\newtheorem{theorem}{Theorem}[section]

% share theorem counter with these (they will be Theorem 2.1, 2.2, ...)
\newtheorem{lemma}[theorem]{Lemma}
\newtheorem{proposition}[theorem]{Proposition}
\newtheorem{corollary}[theorem]{Corollary}
\newtheorem{rem}[theorem]{Remark}        % remarks share the sequence
\newtheorem{example}[theorem]{Example}  % examples share the sequence

% definitions use their own counter (won't affect theorem/lemma numbering)
\newtheorem{defn}{Definition}[section]

% unnumbered environments kept as-is
\newtheorem*{pf}{Proof}
\newtheorem*{mthm}{Main Theorem}
\newtheorem*{conjecture}{Conjecture}
\newenvironment{proofthm}[1]{%
  \begin{proof}[Proof of Theorem~#1]%
}{%
  \end{proof}%
}

\newcommand{\norm}[1]{\left\lVert#1\right\rVert}
\newcommand{\M}{\mathcal{M}(n,\xi,r_0)}
\begin{document}
\title{A Cheng-type Eigenvalue-Comparison Theorem for the Hodge Laplacian }
\author[A. Bhattacharya and S. Maity]{Anusha Bhattacharya and Soma Maity}

\address{Department of Mathematical Sciences, Indian Institute of Science Education and Research Mohali, \newline Sector 81, SAS Nagar, Punjab- 140306, India.}
\email{ph21059@iisermohali.ac.in}


\address{Department of Mathematical Sciences, Indian Institute of Science Education and Research Mohali, \newline Sector 81, SAS Nagar, Punjab- 140306, India.}
\email{somamaity@iisermohali.ac.in}


\begin{abstract}
We consider the class of closed Riemannian $n$-manifolds with Ricci curvature and injectivity radius bounded below by uniform constants, and an upper bound on the diameter. We establish a uniform upper bound for the eigenvalues of the Hodge Laplacian acting on differential forms on Riemannian manifolds in this class, similar to the classical eigenvalue comparison theorem proved by Cheng for the Laplace-Beltrami operator acting on smooth functions. This extends earlier work of  Dodziuk and Lott, which required sectional curvature bounds in addition to bounds on other geometric quantities. As an application, we obtain uniform eigenvalue estimates for the connection Laplacian acting on $1$-forms.

\end{abstract}




\maketitle
\section{Introduction} Consider a closed Riemannian manifold $(M,g)$. Let $\lambda_k$  denote the $k$-th positive eigenvalue of the Laplacian (non-negative) acting on smooth functions on $M$, counted with multiplicity. In 1975, S.-Y. Cheng proved the following upper bound for $\lambda_k$.
\begin{theorem}\cite{Cheng} Suppose $M$ is a closed Riemannian manifold with dimension $n$ and diameter $D$. If Ricci curvature of $M$ is greater than or equal to $(n-1)\xi$ then
$$\lambda_k(M)\leq \lambda_0^D(B_\xi(\frac{D}{2k}))$$
where $\lambda_0^D(B_{\xi}\left(r)\right)$ denotes the first positive Dirichlet eigenvalue of the Laplacian on a ball of radius $r$ in the space of constant curvature $\xi$ and dimension $n$. 
\end{theorem}
  For  $1\le p\le n$, the Hodge Laplacian $\Delta$ acting on differential forms of degree $p$ is defined as 
\[\Delta=dd^*+d^*d\] where $d$ is the exterior derivative and  $d^*=(-1)^{n(p+1)+1}*d*$, where $*$ is the Hodge operator. Let $\lambda_{k,p}(M)$ denote the $k$-th eigenvalue of $\Delta$ acting on differential forms of degree $p$.  Upper bounds for  $\lambda_{k,p}(M)$ are obtained by Dodziuk in \cite{ dodziuk1982} and Lott in \cite{LottJ} under sectional curvature bounds extending Cheng's result to differential forms. It is natural to ask whether similar bounds can be established under weaker curvature conditions. In this paper, we obtain a Cheng-type upper bound for manifolds with  Ricci curvature bounded below. 
\begin{theorem} \label{thm:main_theorem}
Let $(M,g)$ be a closed Riemannian $n$-manifold with diameter $D$ such that the Ricci curvature satisfies $\mathrm{Ric}_g \ge (n-1)\xi .$  If the harmonic radius of $(M,g)$ is greater than $r_H$ then 
\[\lambda_{k,p}(M)\le 2^{2p+1} \lambda_0^D\left(B_\xi\left(\frac{D}{2k}\right)\right), \quad \forall k\ge \frac{D}{2r_H}\]
and \[\lambda_{k,p}(M)\le 2^{2p+1} \lambda_0^D\left(B_\xi(r_H)\right),\qquad \forall k\le \frac{D}{2r_H}\]
where $\lambda_0^D(B_{\xi}\left(r)\right)$ denotes the first positive Dirichlet eigenvalue of the Laplacian acting on functions on a ball of radius $r$ in the space of constant curvature $\xi$ and dimension $n.$
\end{theorem}

For the definition of harmonic radius, we refer to section \ref{prelim}. In \cite{lott2022}, Lott mentioned the existence of uniform bounds on eigenvalues under the assumptions mentioned in Theorem \ref{thm:main_theorem}. We remark that $\lambda_{1,p}(M)$ satisfies the upper bound given in the second inequality of the above theorem. Using the estimates for $\lambda_0^D(B_{\xi}\left(r)\right)$ given in \cite{Cheng}, we obtain explicit bounds for $\lambda_{k,p}(M)$ in terms of $n,D,p,\xi$ and $k.$

In a similar direction, Savo derived upper bounds for the spectral gap of the Hodge Laplacian on the hyperbolic space and established bounds for the first eigenvalue of the Hodge Laplacian on convex Euclidean domains with absolute and relative boundary conditions in \cite{savo09} and \cite{savo2002}, respectively. In a subsequent work \cite{guerinisavo2004}, Guerini and Savo derived an estimate of the first eigenvalue of $\Delta$ for manifolds whose boundaries have some degree of convexity.   In \cite{takahashi2001}, \cite{takahashi2003}, Takahashi derived bounds on the spectral gap between the first eigenvalues of the Laplacian on functions and on $1$-forms.


\par \textbf{Idea of the proof:}  Our approach is based on a discretization of the manifold by small geodesic balls.  Several works, including \cite{ burago, belkin, Tatiana}, and more recently \cite{MaityBhattacharya2026} have demonstrated that discretizing a Riemannian manifold provides an effective method for studying the spectrum of the Laplacian.   Since the only curvature assumption imposed is a lower bound on the Ricci curvature, our methods mostly rely on the existence of harmonic coordinate charts and a uniform harmonic radius. We consider disjoint balls of radius smaller than the harmonic radius centered at the points of the discretization of the manifold and use bounds on harmonic coordinates in  \cite{CheegerAnderson} to derive an upper bound for the Dirichlet eigenvalues on the balls. In Theorem \ref{thm:main_theorem}, we then establish an upper bound for the eigenvalues on manifolds of diameter $D$ by considering a geodesic of maximal length, partitioning it into segments of length controlled by the harmonic radius, and discretizing the manifold by extending the balls arising from this partition. We derive two immediate corollaries of Theorem \ref{thm:main_theorem}  according to the sign of the Ricci curvature lower bound and obtain an upper bound in terms of the volume in Theorem \ref{thm:volume}. In Corollary \ref{cor:connectionLap}, using the B\"ochner (Weitzenb\"ock) formula for non-negative Ricci curvature, we derive bounds for the first eigenvalue of the connection Laplacians on $1$-forms. 
On non-compact manifolds, we derive an upper bound for the spectral gap of the Hodge-Laplacian by considering an exhaustion of the manifold using compact balls. We apply Theorem \ref{thm:main_theorem} to these balls to derive an upper bound of the first Dirichlet eigenvalue and take their limit to obtain an upper bound on the spectral gap.
 

\section{Preliminary results} \label{prelim}
In this section, we introduce the definitions and results we will use in this paper. For fixed $n\ge 2$, $\xi\ge 0$, and $r_0>0$, we define
\[
\mathcal{M}(n,\xi,r_0)
=\{(M,g)\;\text{is a complete Riemannian $n$-manifold}:  \mathrm{Ric}_g \ge (n-1)\xi,\; \mathrm{Inj}(M,g)\ge r_0\}.
\]
\subsection{Harmonic Coordinates} \label{sec:TheClass}

Anderson proved that if a 
Riemannian manifold $(M,g)$ belongs to $\mathcal{M}$, then a harmonic coordinate chart of uniform radius exists. This will serve as the basic analytic framework for all local estimates in this article. We state it below.

\begin{theorem}{\cite{CheegerAnderson,Hebey}}
\label{thm:harmonic_radius}
Let $(M,g)\in \mathcal{M}(n,\xi,r_0)$. For $n<q<\infty$, there exists a positive constant $r_H=r_H(q, (n,\xi,r_0))<r_0$ such that on any ball of radius $r_H$  a harmonic coordinate chart $\{x_i\}_{i=1}^n$ is defined. If $g_{ij}=g\big{(}\frac{\partial}{\partial x_i},\frac{\partial}{\partial x_j}\big{)}$  then
\begin{enumerate}
    \item $\frac{1}{2}\delta_{ij}\leq g_{ij}\leq 2\delta_{ij}$, and
    \item $r_H^{1-n/q}\|\partial g_{ij}\|_{L^q}\leq 1$.
\end{enumerate}
\end{theorem}

From here onwards, $r_H$ denotes the uniform harmonic radius obtained from
Theorem~\ref{thm:harmonic_radius}.
\begin{defn}[$\varepsilon$-discretization] \label{defn:epsilondiscretization} For any positive $\epsilon$, a finite set $X\subset M$ is called an \emph{$\varepsilon$--discretization} of $(M,g)$ if:
\begin{enumerate}
    \item $d(p,q)\ge 2\varepsilon$ for all distinct $p,q\in X$;
    \item the family of balls $\{B(p,2\varepsilon)\}_{p\in X}$ covers $M$.
\end{enumerate}

\end{defn}

\par We state a standard comparison principle for eigenvalues of quadratic forms under injective linear maps.
\begin{lemma}
\label{lem:comparing eigenvalues by injective map}
Let $(H_1,\langle\cdot,\cdot\rangle_1)$ and $(H_2,\langle\cdot,\cdot\rangle_2)$ be Hilbert spaces, and let
\[
\Phi : H_1 \to H_2
\]
be a linear injective map.  
Let $Q_1$ and $Q_2$ be quadratic forms on $H_1$ and $H_2$, respectively.  
Assume that there exist positive constants $C_1,C_2>0$ such that for all $f\in H_1$,
\[
\langle f,f\rangle_1 \le C_1 \langle \Phi(f),\Phi(f)\rangle_2,
\qquad
Q_1(f) \ge C_2\, Q_2(\Phi(f)).
\]
Then, for every $k\in\mathbb{N}$,
\[
\lambda_k(Q_1) \ge \frac{C_2}{C_1}\,\lambda_k(Q_2).
\]
\end{lemma}

The next lemma provides a domain decomposition estimate for the Hodge Laplacian. This allows us to control global eigenvalues in terms of local Dirichlet problems.

\begin{lemma}
\label{lem: manifoldanddirichlet}    
Let $(M,g)$ be a closed Riemannian $n$-manifold and $D_1,\dots,D_j \subset M$ be pairwise disjoint domains in M. Let $\lambda^{D}_{j,p}(D_i)$ be the eigenvalue of the Hodge Laplacian $\Delta$ acting on $p$-forms on $D_i$
with Dirichlet boundary conditions. Then, 
\begin{equation*}
\lambda_{jl - 1,p}
\;\le\;
\max_{1 \le i \le j} \lambda^{D}_{l-1,p}(D_i). 
\end{equation*}
\end{lemma}

\begin{proof}

Define the Hilbert space $H_1 = \bigoplus_{i=1}^{j} H^1_0\!\left(\Lambda^p T^*D_i\right)$ and for $\omega=\left(\omega_1,\omega_2,\dots,\omega_j\right)$ let
\[
\langle \omega,\omega \rangle_{H_1}
=
\sum_{i=1}^{j} \int_{D_i} |\omega_i|^2 \, dv_g,
\]
and quadratic form
\[
Q_1(\omega)
=
\sum_{i=1}^{j}
\int_{D_i}
\bigl(
|d\omega_i|^2 + |d^*\omega_i|^2
\bigr)\, dv_g.
\]


Each $\omega_i$ is extended by zero outside $D_i$, which is admissible since the domains are pairwise disjoint and the boundary conditions are Dirichlet.
From Lemma \ref{lem:comparing eigenvalues by injective map},
\[
\frac{Q_1(\omega)}{\langle \omega,\omega \rangle_{H_1}}
\;\le\;
\max_{1 \le i \le j} \lambda^{D}_{l-1,p}(D_i).
\]
Hence the min-max principle for $\Delta$ on $M$ yields the result.
\end{proof}

As a consequence of the domain decomposition estimate, we obtain an upper bound on the global eigenvalues in terms of local Dirichlet problems on the discretization balls.

\begin{corollary}
    \label{cor:upper-bound-max}
Let $\varepsilon<r_H$, $X=\{x_i\}$ be an $\varepsilon$--discretization of $M$ and $\{B_i=B(x_i,\varepsilon)\}$.  Then,
\[\lambda_{|X|-1,p}\leq\max_i\lambda_{0,p}^{D}(B_i).\]
\end{corollary}


%%%%%%%%%%%%%%%%%%%%%%%%%%%%%%%%%%%%%%%%%%%%%%%%%%%%%%%%%%%%%%%%%%%%%%
\section{Upper Bound of the Hodge Laplacian Eigenvalues}\label{sec:upperbound}
In this section, we derive an upper bound for the eigenvalues of the Hodge Laplacian. We combine local Dirichlet eigenvalue estimates with a domain decomposition to derive a global, Cheng-type estimate.

The following lemma establishes a uniform upper bound for the Dirichlet eigenvalues on geodesic balls of radius $\varepsilon$ smaller than the harmonic radius $r_H$ for manifolds  $(M,g) \in \mathcal{M}(n,\xi,r_0)$. 
This estimate will later be used to derive upper bounds for $\lambda_{k,p}$ on balls centered at the points of an $\varepsilon$-discretization of $M$.
\begin{lemma}    \label{lem:upper bound dirichlet eigenvalue}
Let $(M,g)\in \mathcal{M}(n,\xi,r_0)$ and $x\in M.$ For $0<\varepsilon\le r_H$,  let $B=B(x, \varepsilon)$ and $B_\xi(\varepsilon)$ be a ball of radius $\varepsilon$ in the space of constant curvature $\xi.$ Then   
\[\lambda_{0,p}^D(B)\le 2^{2p+1}\lambda_0^D(B_\xi(\varepsilon))\]
where $\lambda_0^D(B)$ denotes the first Dirichlet eigenvalue of the Laplace-Beltrami operator acting on functions.
\end{lemma}
\begin{proof}
Since $\varepsilon<r_H$, we choose  harmonic coordinates $\{y_j\}_{j=1}^n$ in $B$ and  define
\[
\omega_0 = dy^{1}\wedge\cdots\wedge dy^{p}.
\]
For a $p$ form $\alpha=\sum_{i_1<i_2<\cdots <i_{p}}\alpha_{i_1i_2\cdots  i_{p}}dx^{i_1}\wedge dx^{i_2}\cdots dx^{i_{p}}$, from (3.3.47) \cite{jost}
\begin{equation}
\label{codifferential_expansion}
    \left(d^*\alpha\right)_{i_1i_2\cdots i_{p-1}}=-g^{kl}\left(\frac{\partial \alpha_{ki_1i_2\cdots i_{p-1}}}{\partial x^l}-\Gamma^j_{kl}\alpha_{ji_1\cdots i_{p-1}}\right).
\end{equation}
For harmonic coordinates $\{y_j\}_{j=1}^n$ in $B$,  we have $g^{kl}\Gamma_{kl}^j=0$ for all $j=1,2,\cdots,n$. Hence, from (\ref{codifferential_expansion}) 
\begin{equation}
    \label{codifferential of omega0}
    d^*\omega_0=0.
\end{equation}

Let $f\in H^1_0(B)$ and $\omega=f\omega_0$. By the minmax principle,
\begin{equation}
    \label{eq:rayleigh_quotient_dirichlet}
    \lambda_{0,p}^{D}(B)\le \frac{\int_{B} (|d\omega|^2+|d^*\omega|^2)v_g}{\int_{B}|\omega|^2 v_g}.
\end{equation}

Now, \begin{equation} \label{eq:domegabound}
 |d\omega|^2=|df\wedge\omega_0|^2\le |df|^2|\omega_0|^2
\end{equation}
On the other hand, using (\ref{codifferential of omega0}), 
\begin{equation}
\label{eq:codifferential_in_coords}
    d^*f\omega_0= \left(g^{1l}\frac{\partial f}{dx^l}\right)dx^2\wedge \cdots\wedge dx^p.
\end{equation}
Using diagonalization and change of variables in Theorem \ref{thm:harmonic_radius} we can derive for all $1\le i,j\le n$,
\begin{equation}
 \frac{1}{2}\delta_{ij}\le g^{ij}\le 2\delta_{ij} \nonumber
\end{equation}
where $(g^{ij})=(g_{ij})^{-1}.$ These also give bounds on the eigenvalues of all submatrices of $(g^{ij})$ and hence, their determinants. Also $|dx^1(\triangledown f)|^2\le |df|^2 .$
Using these bounds in (\ref{eq:codifferential_in_coords}), we have
\begin{equation}
    \label{eq:norm_codifferential_bound}
    |d^*f\omega_0|^2\le 2^{p}|df|^2.
\end{equation} 
Also 
\begin{equation}
  \label{eq:bound_on_norm_omega0}  
  |\omega_0|^2\ge \frac{1}{2^p}.
\end{equation}

Substituting (\ref{eq:domegabound}), (\ref{eq:norm_codifferential_bound}) and (\ref{eq:bound_on_norm_omega0}) in (\ref{eq:rayleigh_quotient_dirichlet}),
\[\lambda_{0,p}^D(B)\le 2^{2p+1} \frac{\int_B |df|^2v_g}{\int_B|f|^2v_g}\]
for all $f\in H^1_0(B)$. Hence, 
\[\lambda_{0,p}^D(B)\le 2^{2p+1} \lambda_{0,p}^D(B).\] 

Hence, from Theorem 1.1 in \cite{Cheng},
\[\lambda_{0,p}^D(B)\le 2^{2p+1}\lambda_0^D(B_\xi(\varepsilon)).\]

\end{proof}

In particular, we obtain a similar result related to Lott's eigenvalue bounds \cite{LottJ} where we use Ricci curvature and injectivity radius bounds instead of sectional curvature and volume constraints. 
\subsection{Proof of Theorem \ref{thm:main_theorem}:}

\begin{proof}
    Since $(M,g)$ is closed, by the Hopf-Rinow theorem, there exists a minimizing geodesic $\gamma$ of length $D$. We choose the following decomposition of $\gamma$:
    \[\gamma(0),\gamma\left(\frac{D}{k}\right),\gamma\left(\frac{2D}{k}\right),\cdots,\gamma\left(\frac{(k-1)D}{k}\right),\gamma(D).\]
Since $\gamma$ is minimizing, each of the points in the above partition are equally spaced by distance $\frac{D}{k}$. Let $B_i=B\left(\gamma\left(\frac{iD}{k}\right),\frac{D}{2k}\right)$. For $i\neq j$, if $p\in B_i\cap B_j$, then \[d\left(\gamma\left(\frac{iD}{k}\right),\left(\frac{jD}{k}\right)\right)\le d\left(p,\left(\frac{iD}{k}\right)\right)+ d\left(p,\left(\frac{jD}{k}\right)\right)<\frac{D}{k}\] which is a contradiction. Hence, $B_i\cap B_j=\phi$ for all $i\neq j.$ \\ Choose $\varepsilon= \frac{D}{2k}$ and let $X(\varepsilon)$ be an $\varepsilon$-discretization of $M$ by choosing an extension of the balls $B_i$ defined above. Clearly, $|X(\varepsilon)|\ge k+1$. For $k\ge\frac{D}{2r_H}$, we have $\varepsilon\le r_H$. 

Hence, we can apply Lemma \ref{lem:upper bound dirichlet eigenvalue}
and conclude that 
\[\lambda_{k,p}\le 2^{2p+1} \lambda_0^D\left(B_\xi\left(\frac{D}{2k}\right)\right), \quad \forall k\ge \frac{D}{2r_H}\]
and \[\lambda_{k,p}\le 2^{2p+1} \lambda_0^D\left(B_n(\xi,r_H\right)),\qquad \forall k\le \frac{D}{2r_H}.\]
\end{proof}

As a consequence of the above theorem, we obtain the following corollaries.

\begin{corollary}
\label{cor:positive_ricci}
Let $(M,g)$ be a closed Riemannian manifold with non-negative Ricci curvature.  If the harmonic radius is bounded below by $r_H$ and $\mathrm{ diam}(M)= D$ then 

\[\lambda_{k,p}\le 2^{2p+1}n^2\pi^2\frac{k^2}{D^2},\qquad \forall k\ge\frac{D}{2r_H}\]
and \[\lambda_{k,p}\le 2^{2p-1}n^2\pi^2r_H^{-2},\qquad \forall k\le \frac{D}{2r_H}.\]

\end{corollary}

\begin{proof}
    From the proof of Theorem 2.1 in \cite{Cheng}, we observe that there is an n-cube with sides $\frac{2}{\sqrt{n}}\frac{D}{2k}$ contained in $B_n(\xi,\frac{D}{2k})$. Using the first eigenvalue of the cube and the domain monotonicity theorem, the result follows.
\end{proof}
For a non-positive lower bound on the Ricci curvature, we have the following upper bound on $\lambda_{k,p}.$ 
\begin{corollary} Let $(M,g)$ be a closed Riemannian manifold with $\rm{Ric}\geq -(n-1)\xi$.  If the harmonic radius is bounded below by $r_H$ and $\mathrm{diam}(M)= D$ then 
\begin{enumerate}
    \item For $n=2(m+1)$ where $m=0,1,2,\cdots$

\[\lambda_{k,p}\le 2^{2p-1}(2m+1)^2\xi+2^{2p+3}(1+2^m)^2\pi^2\frac{k^2}{D^2},\qquad \forall k\ge\frac{D}{2r_H}\]
and \[\lambda_{k,p}\le 2^{2p-1}(2m+1)^2\xi+2^{2p+1}(1+2^m)^2\pi^2r_H^{-2},\qquad \forall k\le \frac{D}{2r_H}.\]
\item For $n=2m+3$, $m=0,1,2,\cdots$
                    \[\lambda_{k,p}\le 2^{2p-1}(2m+2)^2\xi+2^{2p+3}(1+2^{2m})^2(1+\pi^2)\frac{k^2}{D^2},\qquad \forall k\ge\frac{D}{2r_H}\]
     and \[\lambda_{k,p}\le 2^{2p-1}(2m+2)^2\xi+2^{2p+1}(1+2^{2m})^2(1+\pi^2)r_H^{-2},\qquad \forall k\le \frac{D}{2r_H}.\]       
\end{enumerate}  
\end{corollary}
\begin{proof}
  Using Corollary 2.3 \cite{Cheng} in Lemma \ref{lem:upper bound dirichlet eigenvalue}, we get the required bounds.
\end{proof}

\begin{theorem}
\label{thm:volume}
 Let $0\le r_H\le \frac{1}{\sqrt{|\xi|}}$ and $(M,g)$ be a closed Riemannian $n$-manifold with volume $V$. Assume that the Ricci curvature satisfies $\mathrm{Ric}_g \ge (n-1)\xi .$ Then   for $n>1$,

\[
\lambda_{k,p}
\le
2^{2p+n+5}\cdot (k+1)\cdot \left(\frac{\alpha_n}{V}\right)^\frac{2}{n},\qquad \mathrm{for }\hspace{1mm} k>\frac{2^nV}{\alpha_{n-1}r_H^n}.
\]
where $\alpha_n=\mathrm{vol}(S^n).$
\end{theorem}
\begin{proof}
We choose \[\varepsilon
=
\frac{2}{\sqrt{|\xi|}}\operatorname{arcsinh}\left[\left(\frac{ V|\xi|^\frac{n}{2}}{\alpha_n(k+1)}\right)^{1/n}\right].
\]  

From Bishop's theorem, 
\begin{equation}
\label{eq:volume_ratio}
  |X(\varepsilon)|\ge \frac{V}{\sup_{x\in X(\varepsilon)}\{B(x,\varepsilon/2)\}}\ge \frac{V}{V_{\xi}(\varepsilon/2)}.  
\end{equation}
\begin{align}
\mathrm{Vol}_{\xi}(\varepsilon/2)
&=
\alpha_{n-1}
\int_{0}^{\varepsilon/2}
\left(
\frac{\sinh(\sqrt{|\xi|}t)}{\sqrt{|\xi|}}
\right)^{n-1}
dt \nonumber\\
&=
\frac{\alpha_{n-1}}{|\xi|^{n/2}}
\int_0^{\varepsilon \sqrt{|\xi|}/2}
\sinh^{n-1}(t)\,dt \nonumber\\
&\le
\frac{\alpha_{n-1}}{|\xi|^{n/2}}
\left(
\sinh\left(\frac{\varepsilon \sqrt{|\xi|}}{2}\right)
\right)^n \nonumber
\end{align}

Substituting the above bound and chosen $\varepsilon$ in (\ref{eq:volume_ratio}),

\begin{align}
|X(\varepsilon)|\ge 
\frac{V |\xi|^{n/2}}
{\alpha_{n-1}
\left(
\sinh\left(\frac{\varepsilon \sqrt{|\xi|}}{2}\right)
\right)^n} \ge k+1.
\end{align}
Now $y\le \sinh{y}\implies \frac{1}{y^n}\ge \frac{1}{(\sinh y)^n}$. Hence for $k>\frac{2^nV}{\alpha_{n-1}r_H^n}$, we have $k>\frac{V|\xi|^\frac{n}{2}}{\alpha_{n-1}(\sinh(r_H\sqrt{|\xi|}/2))^n}.$ Hence,  $\varepsilon<r_H$ for $k>\frac{2^nV}{\alpha_{n-1}r_H^n}$ . Using Lemma \ref{lem:upper bound dirichlet eigenvalue}, for $f:B(x,\varepsilon)\rightarrow\mathbb{R}$ defined as $f(y)=1-\frac{1}{\varepsilon}d(x,y)$. Then $|df|\ge \frac{1}{\varepsilon}$ and $|f|\ge\frac{1}{2}$ in $B\left(x,\frac{\varepsilon}{2}\right)$. Hence, using $\omega=f\omega_0$ as the test form in the lemma and using Bishop-Gromov volume ratio, we have
\begin{equation}
    \label{eq:lambda upper bound for volume}
    \lambda_{k,p}\le 2^{2p+3}\frac{V_\xi(\varepsilon)}{\varepsilon^2 V_\xi(\frac{\varepsilon}{2})}, \qquad \mathrm{for}\hspace{1mm} k>\frac{2^nV}{\alpha_{n-1}r_H^n}.
\end{equation}
Let $x^n=\frac{ V|\xi|^\frac{n}{2}}{\alpha_n(k+1)}$. Then \begin{align}
    \frac{V_\xi(\varepsilon)}{ V_\xi(\frac{\varepsilon}{2})} &\le \left[\frac{\sinh\left(\varepsilon\sqrt{|\xi|}\right)}{\sinh\left(\varepsilon\sqrt{|\xi|}/2\right)}\right]^n \nonumber\\
    &\le 2^n\left[\cosh\left(\varepsilon\sqrt{|\xi|}/2\right)
    \right]^n \nonumber \\
    &\le 2^n\left[\cosh\left(\sinh^{-1}x\right)\right]^n \nonumber\\
    &\le 2^n\left(1+x\right)^n.
\end{align}
Since $\sinh^{-1}x\ge \frac{x}{1+x}$, substituting $\varepsilon$ in (\ref{eq:lambda upper bound for volume})  gives
\begin{equation}
\label{eq:upper bound after sinh substitution}
    \lambda_{k,p}\le 2^{2p+3}|\xi|\frac{(1+x)^{n+2}}{x^2}.
\end{equation}
Now $r_H\le \frac{1}{\sqrt{|\xi|}}$implies $\sinh \left(r_H\sqrt{|\xi|}/2\right)<1$. Hence,  $k>\frac{2^nV}{\alpha_{n-1}r_H^n}>\frac{V|\xi|^\frac{n}{2}}{\alpha_{n-1}(\sinh(r_H\sqrt{|\xi|}/2))^n}> \frac{V|\xi|^\frac{n}{2}}{\alpha_{n-1}}$ which implies $x<1.$\\

Substituting in (\ref{eq:upper bound after sinh substitution}),
\[ \lambda_{k,p}\le 2^{2p+n+5}\cdot (k+1)\cdot \left(\frac{\alpha_n}{V}\right)^\frac{2}{n}.\]

\end{proof}
\begin{corollary} \label{cor:connectionLap}
    Let $(M,g)$ be a closed $n$-manifold with ${\rm inj}(M,g) \ge r_0$, ${\rm diam}(M) \le D <\infty$, $ {\rm Ric}_g \ge 0$. Let $r_H$ denote the lower bound on the harmonic radius of M. Then the first non-zero eigenvalue $\lambda^C_{1,1}$ of the connection Laplacian $\nabla^{*} \nabla$ acting on the space of $1$-forms satisfies
   \[ \lambda_{1,1}^C\le \frac{2^{2p+1}n^2\pi^2}{r_H^2}.\]
\end{corollary}
\begin{proof} For a smooth Riemannian manifold $(M,g)$ and a smooth $1$-form $\omega$, we recall the Bochner (Weitzenb\"ock) formula for $1$-forms:
\begin{equation}
\label{eq: weitzenbock pointwise}
\Delta \omega
=
\nabla^{*}\nabla \omega
+
\mathrm{Ric}(\omega),
\end{equation}
where $\nabla^{*}\nabla$ is the rough (connection) Laplacian, and $\mathrm{Ric}(\omega)$ is the action of the Ricci tensor on $\omega$.

Taking the $L^2$-inner product of \eqref{eq: weitzenbock pointwise} with $\omega$  and integrating over $M$, we obtain
\[
\int_M \langle \Delta \omega , \omega \rangle
=
\int_M \langle \nabla^{*}\nabla \omega , \omega \rangle
+
\int_M \mathrm{Ric}(\omega,\omega).
\]

Since $M$ is closed, by integration by parts we obtain
\begin{equation*}
\int_M \bigl( |d \omega|^2 +  |\delta \omega|^2\bigr)
=
\int_M |\nabla \omega|^2
+
\int_M \mathrm{Ric}(\omega,\omega).
\end{equation*}

For $\mathrm{Ric}\ge 0$, on $1$-forms we have $\mathrm{ker}(\nabla^*\nabla)=\mathrm{ker}(\Delta)$. Also,
\begin{equation*}
   \frac{\int_M |\nabla \omega|^2}{\int_M |\omega|^2}\le \frac{\int_M \bigl( |d \omega|^2 +  |\delta \omega|^2\bigr)}{\int_M | \omega|^2 }.
\end{equation*}
Applying min-max principle,
\[\lambda_{1,1}^C\le \lambda_{1,1}(M).\]
We then use Corollary \ref{cor:positive_ricci} and domain monotonicity of eigenvalues to obtain the desired inequality.
\end{proof}

\section{Bound on spectral gap for non-compact manifolds} \label{sec:non_compact}
In this section, we study bounds for the spectral gap of the Hodge Laplacian acting on $p$-forms on complete non-compact Riemannian manifolds. Since the spectrum may contain continuous components in the non-compact setting, we work with the supremum of the $L^2$ spectrum. Using harmonic coordinate estimates in Theorem \ref{thm:harmonic_radius} and comparing with Dirichlet eigenvalues on a sequence of compact balls with increasing radius, we derive upper bounds for the spectral gap under Ricci curvature lower bounds. 
\begin{defn}
    Let M be a $n$-dimensional Riemannian manifold without boundary. The  supremum of $L^2$ spectrum of $p$-forms is defined as
    \[\sigma^p(M)=\inf_{\omega\in \Lambda^p_0(M)}\frac{\int\|d\omega\|^2+\|d^*\omega\|^2}{\int \|\omega\|^2}.\]
    
\end{defn}
The following corollary follows from Lemma \ref{lem:upper bound dirichlet eigenvalue}.
\begin{corollary} For $r_H\rightarrow \infty$ and $\mathrm{Ric}(M)\ge -(n-1)\xi$, 
    \[\sigma^p(M)\le 2^{2p-1}(n-1)^2\xi.\]
\end{corollary}
\begin{proof}
    For $r_H \rightarrow \infty$, we have a harmonic coordinate chart on the entire manifold along with bounds on metric coefficients given in Theorem \ref{thm:harmonic_radius}. Hence, following the proof of Lemma \ref{lem:upper bound dirichlet eigenvalue}, we have  \[\sigma^p(M)\le 2^{2p+1} \sigma^0(M)\]
    where $\sigma^0(M)=\inf_{f\in C_0^\infty(M)}\frac{\int_M |df|^2}{\int_M|f|^2}.$ From Theorem 4.2 \cite{Cheng}, we obtain the result.
    
    \end{proof}

We have the following estimate of the supremum of $L^2$ spectrum of $p$ forms using the Dirichlet eigenvalues on  an exhaustion of the manifold by compact submanifolds 

\begin{lemma}
\label{lem:exhausting_M_dirichlet}
 Let $M$ be a Riemannian $n$-manifold and $M_i\subseteq M$ such that $M_i$ is compact for each $i\in \mathbb{N}.$  If  $\bigcup_i M_i=M$, then
    \[\sigma^p(M)\le\lim_{i\rightarrow\infty} \lambda_{0,p}^D(M_i).\]
\end{lemma}

\begin{proof}
    Since $\Lambda^p_0(M_i)\subseteq \Lambda^p_0(M_{i+1})$, the lemma follows.
\end{proof}

Savo \cite{savo09}  obtained the following bound on the supremum of $L^2$ spectrum on $p$ forms on the hyperbolic space $\mathbb{H}^n$.
\begin{equation*}
\sigma^p(\mathbb{H}^n)=
\begin{cases}
0, & \text{if $p\le (n+1)/2$, }\\
\frac{(2p-n-1)^2}{4}, & \text{if $p>(n+1)/2$}.
\end{cases}
\end{equation*}
Savo also showed that if  $M$ admits a non-trivial harmonic $p$-form then $\sigma_p(M)=0$ \cite{savo09}. We now derive a general upper bound for the supremum of the $L^2$ spectrum on complete non-compact manifolds whose Ricci curvature admits a uniform lower bound.
 
\begin{theorem}
   
   Let $(M,g)$ be a complete non-compact Riemannian $n$-manifold such that $\mathrm{Ric}(M)\ge (n-1)\xi.$ Then there exists a positive constant $r_H$ such that for $p$ forms
   \[\sigma^p(M)\le 2^{2p+1} \lambda_0^D(B_\xi(r_H)).\]   
   \end{theorem}
  \begin{proof}  
  

 Let $x\in M$ and $\{B_i=\overline{B\left(x,i\right)}\}_{i\in\mathbb{N}}$ be an exhaustion of $M$ by compact domains. Fixing a lower bound $r_0$ of the injectivity radius, from \cite{Hebey}, the harmonic radius is bounded below by a constant $r_H$.  We choose $j\in\mathbb{N}$ such that  $2r_H\le j.$      Then for all $j'>j$, from Theorem \ref{thm:main_theorem}, 
\[\lambda_{0,p}^D(B_j')\le 2^{2p+1} \lambda_0^D(B_\xi(r_H)).\]      
     Since $\sigma^p(M)\le \lim_{i\rightarrow \infty}\lambda_{0,p}^D(B_i)$, \[\sigma^p(M)\le 2^{2p+1} \lambda_0^D(B_\xi(r_H)).\]

  \end{proof}      





\section*{Acknowledgement}
The authors are grateful to Professor John Lott and Professor Shouhei Honda for their helpful comments and suggestions. Anusha Bhattacharya is supported by the Council of Scientific and Industrial Research (CSIR).
\printbibliography
\end{document}
